\scanpage{142}
% archive sheet number: 141

\[
\delta(T)=\delta(T')
\]

donc $T\mapsto T'$ induit un \deleted{\unsure{$k$}} autom
de $\widetilde S=X\wedge X'=\delta^{-1}(\unsure{-|{}}\text{section nulle de }\unsure{\check V|V^\vee|\mathbb V}(L\otimes L'))$.

Quand on choisit bases $u,v$ donc
$L_P$ identifié à l'axe des $t$, l'involution
$T\mapsto T'$ \deleted{à qui} se représente par

\[
t\longmapsto bb'-t,
\]

qui \deleted{\unsure{est|$\Rightarrow$}} \inserted{induit} l'\uline{involution can.} sur $\widetilde S$, schéma
des solutions de

\[
t^2-bb't+\text{fourbis}=0.
\]

\deleted{\unsure{A|$\triangle$}} On va donner une interprétation
schématique de $X\wedge X'$ comme produit
contracté.

\uline{Prop} \deleted{\unsure{Soit}}Supposons que 2 soit inv., ou
plus généralement, qu'en tt pt où 2 est nul,
soit $X$ soit $X'$ est étale (cela signifie
aussi que $2,\underset{\substack{\parallel\\ b^2-4c}}{\deleted{\unsure{b}}\delta},\underset{\substack{\parallel\\ b'{}^2-4c'}}{\deleted{\unsure{b'}}\delta'}$ ne s'annulent pas
simultanément sur $S$, i.e que
l'idéal $(2,\delta,\delta')$ de $k$ est l'idéal unité).
Alors le \deleted{\unsure{hom}} morphisme

\diagram{\begin{tikzcd}[ampersand replacement=\&]
X\times_S X'\arrow[r] \& X\wedge_S X'
\end{tikzcd}}

est \deleted{\unsure{içi fid|fid}} un morphisme quadratique
(fini loc\textsuperscript{t} libre de rang 2), et si $\sigma_X,\sigma_{X'}$ sont
les inv can. de $X$ et $X'$ resp\textsuperscript{t}, $\sigma_X\times_S\sigma_{X'}$ est l'inv. can

